\documentclass[12pt,a4paper]{article}

\usepackage[spanish]{babel}
\usepackage[utf8]{inputenc}
\usepackage[T1]{fontenc}
\usepackage{geometry}
\usepackage{graphicx}
\usepackage{amsmath}
\usepackage{booktabs}
\usepackage{longtable}
\usepackage{hyperref}
\usepackage{enumitem}
\usepackage{titlesec}
\usepackage{amstext}

\geometry{margin=2.5cm}

\title{\textbf{Hito 1: Alcance y viabilidad}}
\author{
Asignatura: Desarrollo y Despliegue de Soluciones Big Data \\
Máster Universitario en Big Data y Computación en la Nube \\
Curso académico 2025--2026 \\
\\
Alonso Marcos Muñoz \\
\texttt{Alonso.Marcos@alu.uclm.es} \\
\\
Jose Barros Ribademar \\
\texttt{Jose.Barros1@alu.uclm.es}
}
\date{Febrero de 2026}

\begin{document}

\maketitle
\tableofcontents
\newpage

\section{Introducción del hito}

Este primer hito cierra el alcance del proyecto antes de entrar en la fase técnica de implementación. El objetivo no es solo describir el problema de negocio, sino justificar por qué tiene sentido abordarlo con una solución de datos, qué valor económico realista puede aportar y cómo se va a ejecutar con recursos limitados, en el calendario oficial de la asignatura.

La propuesta se ha redactado con un enfoque deliberadamente alcanzable: empezar con una primera versión funcional (MVP) y evitar compromisos técnicos innecesarios para una entrega inicial. Esta decisión mantiene el equilibrio entre rigor académico, viabilidad de ejecución en pareja y potencial de mejora en los hitos siguientes.

Fecha oficial de cierre del hito: \textbf{27 de febrero de 2026}.

\section{Alcance y viabilidad}

\subsection{Definición del problema de negocio}

La operadora de telecomunicaciones gestiona una base de datos activa de 2.000.000 de clientes particulares en un mercado altamente saturado. Actualmente, la compañía se enfrenta a una tasa de rotación mensual del 3\% (60.000 clientes abandonan la compañía cada mes), una hemorragia que las campañas de fidelización genéricas no logran detener.

Esta falta de inteligencia comercial genera un impacto negativo por dos vías:

\begin{itemize}
    \item Incapacidad de anticiparse al descontento: 40.000 usuarios recuperables se pierden cada mes.
    Con un valor de vida promedio perdido de 100 euros de margen neto anual por cliente, la pérdida asciende a:

    \[
    40.000 \times 100 = 4.000.000 \text{ euros/mes}
    \]

    \item Incentivos innecesarios al 5\% de la base leal (100.000 clientes).
    Con un coste medio de 10 euros por descuento erróneo:

    \[
    100.000 \times 10 = 1.000.000 \text{ euros/mes}
    \]
\end{itemize}

El balance total de pérdidas operativas asciende a \textbf{5.000.000 euros al mes}.

\subsection{Planteamiento y selección de la solución técnica}

Se analizaron tres escenarios: (1) heurístico por reglas, (2) machine learning en servidor monolítico y (3) arquitectura big data distribuida.

\begin{longtable}{p{4cm}p{2.5cm}p{2.5cm}p{3cm}}
\toprule
\textbf{Criterio de decisión} & \textbf{Heurística} & \textbf{ML monolítico} & \textbf{Big Data (MVP)} \\
\midrule
Escalabilidad & Baja & Media & Alta \\
Mantenimiento & Alto & Medio & Medio \\
Integración de nuevas fuentes & Baja & Media & Alta \\
Trazabilidad & Baja & Media & Alta \\
Riesgo de obsolescencia & Alto & Medio & Bajo \\
\bottomrule
\end{longtable}

Se selecciona una arquitectura big data distribuida en versión mínima viable (MVP): procesamiento por lotes, primer modelo de clasificación y mejora incremental en hitos posteriores.

\subsection{Evaluación de la viabilidad y valor}

\subsubsection{Viabilidad técnica}

La viabilidad técnica es favorable por cuatro razones:

\begin{enumerate}
    \item Masa crítica suficiente (2 millones de clientes).
    \item Problema de clasificación supervisada natural.
    \item Existencia de señales predictivas razonables.
    \item Diseño batch que reduce riesgo inicial.
\end{enumerate}

\subsubsection{Viabilidad económica (escenario conservador)}

Hipótesis de mejora moderada:

\begin{itemize}
    \item Reducción 4\% fuga recuperable:
    \[
    4.000.000 \times 0.04 = 160.000 \text{ euros/mes}
    \]
    \item Reducción 5\% descuentos innecesarios:
    \[
    1.000.000 \times 0.05 = 50.000 \text{ euros/mes}
    \]
\end{itemize}

Beneficio total estimado:
\[
210.000 \text{ euros/mes}
\]

Costes estimados:
\begin{itemize}
    \item Equipo técnico: 10.000 euros/mes
    \item Infraestructura cloud: 1.100 euros/mes
\end{itemize}

Coste total:
\[
11.100 \text{ euros/mes}
\]

ROI mensual:
\[
\frac{210.000 - 11.100}{11.100} \times 100 = 1.791,89\%
\]

Payback:
\[
\frac{11.100}{210.000} \times 30 = 1,59 \text{ días}
\]

\subsubsection{Viabilidad ética y legal}

Se establecen tres medidas:

\begin{enumerate}
    \item Seudonimización y control de accesos.
    \item Trazabilidad de datasets y modelos.
    \item Métrica de fairness: Equal Opportunity Difference.
\end{enumerate}

\section{Planificación y recursos}

\subsection{KPIs de negocio y traducción técnica}

Objetivos del MVP:

\begin{itemize}
    \item Reducir 4\% fuga recuperable.
    \item Reducir 5\% incentivos innecesarios.
\end{itemize}

Traducción operativa:

\begin{itemize}
    \item Fuga: 40.000 $\rightarrow$ 38.400 clientes/mes.
    \item Incentivos: 100.000 $\rightarrow$ 95.000 clientes/mes.
\end{itemize}

\subsection{Equipo de trabajo}

\begin{itemize}
    \item Ingeniería de datos/MLOps: ingesta, calidad, orquestación.
    \item Modelado: features, entrenamiento y evaluación.
\end{itemize}

\subsection{Fuentes y stack tecnológico}

Fuentes:
\begin{itemize}
    \item Maestro de clientes
    \item Histórico de bajas
    \item Uso y facturación
    \item Histórico de campañas
\end{itemize}

Stack:
\begin{itemize}
    \item Databricks + Delta Lake
    \item Spark (batch)
    \item Spark MLlib
    \item MLflow
    \item GitHub
\end{itemize}

\subsection{Cronograma oficial}

\begin{itemize}
    \item Hito 1: 24–27 febrero 2026
    \item Hito 2: 28 febrero–20 marzo 2026
    \item Hito 3: 21 marzo–17 abril 2026
    \item Hito 4: 18 abril–1 mayo 2026
\end{itemize}

\section{Cierre del hito}

El alcance queda cerrado con una propuesta coherente en las cuatro dimensiones exigidas: definición del problema, selección técnica razonada, viabilidad completa y planificación.

La estrategia adoptada evita promesas difíciles de sostener en febrero de 2026 y prioriza una entrega sólida, medible y defendible. El siguiente paso es ejecutar el hito 2 para validar calidad del dato y materializar la arquitectura de trabajo en Databricks.

\end{document}